\documentclass[../../../include/open-logic-section]{subfiles}

\begin{document}

\olfileid{sth}{ordinals}{opps}
\olsection{الأعداد الترتيبية الخلفية والحدية}

سنستعمل الأعداد الترتيبية كثيرًا في الفصول القليلة التالية. ولذلك من
المفيد أن ندخل بعض المفاهيم البسيطة.

\begin{defn}
لكل عدد ترتيبي~$\alpha$، يكون \emph{خلفه}
$\ordsucc{\alpha} =\alpha \cup \{\alpha\}$. ونقول إن~$\alpha$
عدد ترتيبي \emph{خلفي} إذا كان~$\ordsucc{\beta} = \alpha$ لعدد ترتيبي
ما~$\beta$. ونقول إن~$\alpha$ عدد ترتيبي \emph{حدي} إذا وفقط إذا لم
يكن خاليًا ولا عددًا ترتيبيًا خلفيًا.
\end{defn}
\noindent
تبيّن النتيجة الآتية أن هذا هو المفهوم الصحيح لـ\emph{الخلف}:

\begin{prop}
لكل عدد ترتيبي~$\alpha$:
\begin{enumerate}
	\item $\alpha \in \ordsucc{\alpha}$؛
	\item $\ordsucc{\alpha}$ عدد ترتيبي؛
	\item لا يوجد عدد ترتيبي~$\beta$ يحقق
	$\alpha \in \beta \in \ordsucc{\alpha}$.
	\end{enumerate}
\end{prop}

\begin{proof}
من البديهي أن
$\alpha \in \alpha \cup \{\alpha\} = \ordsucc{\alpha}$. كذلك تكون
$\ordsucc{\alpha}$ مجموعة متعدية من الأعداد الترتيبية، ومن ثم تكون
عددًا ترتيبيًا بموجب~\olref[basic]{corordtransitiveord}.
ويستحيل~$\alpha \in \beta \in \ordsucc{\alpha}$، لأن ذلك يستلزم
$\beta \in \alpha$ أو~$\beta = \alpha$، خلافًا لـ
\olref[basic]{ordordered}.
\end{proof}
\noindent
يسوّغ هذا أيضًا صورة أخرى للبرهان بالاستقراء المتناهي التجاوز:

\begin{thm}[الاستقراء المتناهي التجاوز البسيط]\ollabel{simpletransrecursion}
لتكن~$\phi(x)$ صيغة تحقق:
\begin{enumerate}
	\item $\phi(\emptyset)$؛ و
	\item لكل عدد ترتيبي~$\alpha$، إذا كان~$\phi(\alpha)$، فإن
	$\phi(\ordsucc{\alpha})$؛ و
	\item إذا كان~$\alpha$ عددًا ترتيبيًا حديًا وكان
	$(\forall \beta \in \alpha)\phi(\beta)$، فإن~$\phi(\alpha)$.
\end{enumerate}
عندئذ~$\forall \alpha \phi(\alpha)$.
\end{thm}

\begin{proof}
نثبت المقابل بالنقيض. افترض إذن وجود عدد ترتيبي يحقق~$\lnot\phi$، ولتكن
$\gamma$ أصغر عدد ترتيبي كذلك. فإما أن~$\gamma = \emptyset$، وهذا
يناقض الفرض~1؛ أو أن~$\gamma = \ordsucc{\alpha}$ لبعض~$\alpha$،
وعندئذ تقتضي أصغرية~$\gamma$ أن~$\phi(\alpha)$، فيعطي الفرض~2
$\phi(\gamma)$؛ أو أن~$\gamma$ عدد ترتيبي حدي، وعندئذ تقتضي أصغريته
$(\forall \beta \in \gamma)\phi(\beta)$، فيعطي الفرض~3
$\phi(\gamma)$. وفي الحالات الثلاث نحصل على تناقض.
\end{proof}
\noindent
وسيكون قدر أخير من الترميز نافعًا فيما بعد:

\begin{defn}\ollabel{defsupstrict}
إذا كانت~$X$ مجموعة من الأعداد الترتيبية، فإن
$\supstrict(X) = \bigcup_{\alpha \in X} \ordsucc{\alpha}$.
\end{defn}
\noindent
يرمز الاختصار «lsub» هنا إلى «أصغر حد علوي صارم».
\footnote{تستعمل بعض الكتب «$\text{sup}(X)$» لذلك. لكن كتبًا أخرى
تستعمل «$\text{sup}(X)$» لأصغر حد علوي \emph{غير صارم}، أي ببساطة
$\bigcup X$. وإذا كان لـ$X$ عنصر أكبر~$\alpha$، افترق المفهومان:
فأصغر حد علوي \emph{صارم} هو~$\ordsucc{\alpha}$، في حين أن أصغر حد
علوي \emph{غير صارم} هو~$\alpha$ نفسه.}
وتشرح النتيجة الآتية ذلك:

\begin{prop}
إذا كانت~$X$ مجموعة من الأعداد الترتيبية، فإن~$\supstrict(X)$ أصغر
عدد ترتيبي أكبر من كل عدد ترتيبي في~$X$.
\end{prop}

\begin{proof}
لتكن~$Y = \Setabs{\ordsucc{\alpha}}{\alpha \in X}$، بحيث
$\supstrict(X) = \bigcup Y$. ولما كانت الأعداد الترتيبية متعدية، وكان
كل عنصر من عدد ترتيبي عددًا ترتيبيًا، كانت~$\supstrict(X)$ مجموعة
متعدية من الأعداد الترتيبية، ومن ثم عددًا ترتيبيًا بموجب
\olref[basic]{corordtransitiveord}.

إذا كان~$\alpha \in X$، كان~$\ordsucc{\alpha} \in Y$، ولذلك
$\ordsucc{\alpha} \subseteq \bigcup Y = \supstrict(X)$، ومن ثم
$\alpha \in \supstrict(X)$. وهكذا تكون~$\supstrict(X)$ أكبر تمامًا
من كل عدد ترتيبي في~$X$.

ولتكن~$\gamma$ حدًا علويًا صارمًا اعتباطيًا لـ$X$. إذن
$\alpha \in \gamma$ لكل~$\alpha \in X$. ولما كان~$\gamma$ عددًا
ترتيبيًا، ينتج~$\ordsucc{\alpha} \subseteq \gamma$ لكل
$\alpha \in X$. ومن ثم
$\bigcup Y = \supstrict(X) \subseteq \gamma$، أي
$\supstrict(X) \leq \gamma$. لذلك تكون~$\supstrict(X)$ \emph{أصغر} حد علوي
صارم لـ$X$.
\end{proof}

\end{document}
